% ================================================
% 文件: 03_无源元件.tex
% 章节: 无源元件
% 所属部分: 电子元器件与材料
% 版本: v1.0
% ================================================

\section{无源元件（电阻、电容、电感）}
\label{sec:passive_components}

无源元件是指不需要外部电源就能工作的电子元件，它们不具备放大或有源开关功能。
换言之，无源元件本身既不能向电路净输出能量，也不能通过外加控制量改变信号的功率电平；
它们只能消耗、储存或以无损方式转移电磁能量。这一特征把无源元件与电子管、
晶体管、集成电路等有源器件明确区分开来：有源器件依靠外部直流电源把直流功率
转换为信号功率，从而实现放大与开关；无源元件则只在电压与电流之间建立确定的
约束关系。

从能量的角度看，无源元件可分为两类。第一类是耗能元件，以电阻器为代表，
它把电能不可逆地转换为热能，因而其瞬时功率恒为非负；第二类是储能元件，
以电容器和电感器为代表，它们分别把能量储存在电场和磁场中，并可在随后的
半个周期内把能量归还给电路，因而其瞬时功率可正可负，但在一个完整周期内
平均功率为零。正是由于这种储能与释能的往复过程，电容器与电感器才具有随
频率变化的阻抗，成为滤波、耦合、调谐和阻抗匹配电路的基础。

在通信设备中，无源元件的数量远多于有源器件，其质量与稳定性往往决定整机
的可靠性。发射机的调谐回路、接收机的中频滤波器、电源的整流滤波网络、
天馈系统的匹配电路，无一不是由电阻、电容、电感按特定拓扑组合而成。因此，
维修与设计人员不仅要掌握三种基本元件的理想模型，还必须熟悉其寄生参数、
温度特性与失效模式。理想元件在实际频段中总会伴随寄生量：电阻具有寄生
电感与寄生电容，电容具有等效串联电阻（ESR）与等效串联电感（ESL），
电感具有绕线直流电阻与分布电容。频率越高，这些寄生量的影响越显著。

\begin{table}[htbp]
\centering
\caption{三种基本无源元件的本征关系与能量特性速查}
\begin{tabular}{|p{2.0cm}|p{3.0cm}|p{3.2cm}|p{3.4cm}|}
\hline
\textbf{元件} & \textbf{伏安关系} & \textbf{正弦阻抗} & \textbf{能量特性} \\
\hline
电阻器 $R$ & $u = iR$ & $Z_R = R$（与频率无关） & 耗能，$P = I^2R$ 全部转为热 \\
\hline
电容器 $C$ & $i = C\dfrac{\mathrm{d}u}{\mathrm{d}t}$ & $Z_C = \dfrac{1}{\mathrm{j}\omega C}$ & 储能于电场，$W = \dfrac{1}{2}CU^2$ \\
\hline
电感器 $L$ & $u = L\dfrac{\mathrm{d}i}{\mathrm{d}t}$ & $Z_L = \mathrm{j}\omega L$ & 储能于磁场，$W = \dfrac{1}{2}LI^2$ \\
\hline
\end{tabular}
\end{table}

由表中的阻抗表达式可见，电阻的阻抗与频率无关，电容的阻抗随频率升高而
下降，电感的阻抗随频率升高而上升。这一对偶性质是构造低通、高通、带通
与带阻滤波器的根本依据，也解释了为什么隔直耦合要用电容、而抑制高频干扰
要用电感（磁珠、扼流圈）。

\subsection{电阻器}

电阻器用于限制电流、分压和匹配阻抗。
它是电路中数量最多、结构最简单的元件，其作用可归纳为四类：其一，串联
于回路中限制电流，保护后级器件；其二，与另一电阻构成分压器，为放大级
提供偏置电位；其三，作为负载把有源器件的电流变化转换为电压变化；其四，
作为终端电阻吸收传输线上的入射功率，实现阻抗匹配以消除反射。在射频与
高速数字电路中，第四种用途尤为关键。

电阻器的电阻值取决于导电材料的固有性质与几何尺寸。对于均匀截面的导体，
其电阻为
\[
R = \rho\,\frac{l}{S}
\]
其中 $\rho$ 为材料电阻率，$l$ 为电流通过的长度，$S$ 为截面积。膜式电阻
把高电阻率的碳膜或金属膜沉积在陶瓷基体表面，再用螺旋刻槽的方法延长电流
路径、减小有效截面，从而在很小的体积内获得较大阻值。线绕电阻则把电阻
合金丝绕在绝缘骨架上，阻值精度与功率承受能力都很高，但绕线本身构成螺线管，
带来不可忽略的寄生电感，因此不宜用于高频电路。

电阻器在工作时把电能转换为热能，其消耗的功率为
\[
P = UI = I^2R = \frac{U^2}{R}
\]
这就是电阻的功率计算公式。选用电阻时必须保证实际耗散功率明显低于其
额定功率，工程上通常留有一倍以上的余量，即让实际功率不超过额定值的
$50\%$；在环境温度较高或散热条件不良时还要进一步降额使用。若功率超限，
膜层将因过热而碳化或开裂，表现为阻值漂移、烧焦变色乃至开路。

阻值随温度的变化用温度系数描述。设 $20\,^{\circ}\mathrm{C}$ 时的阻值为
$R_{20}$，温度系数为 $\alpha$，则温度 $t$ 时的阻值近似为
\[
R_t = R_{20}\bigl[\,1 + \alpha (t - 20)\,\bigr]
\]
金属膜电阻的温度系数可低至每摄氏度百万分之几十，碳膜电阻则大得多且多为
负值。除温度系数外，电阻还会产生热噪声（约翰逊噪声），其开路噪声电压
有效值为 $U_n = \sqrt{4kTRB}$，其中 $k$ 为玻尔兹曼常数，$T$ 为绝对温度，
$B$ 为噪声带宽。这说明接收机前端的阻值不宜过大，否则将恶化整机噪声系数。

\begin{itemize}
    \item \textbf{工作原理}：电阻值由材料电阻率、长度和截面积决定：$R = \rho\frac{l}{S}$
    \item \textbf{单位}：欧姆（$\Omega$），常用倍率：k$\Omega$（千欧）、M$\Omega$（兆欧）
    \item \textbf{类型}：碳膜电阻、金属膜电阻、线绕电阻、热敏电阻、压敏电阻
    \item \textbf{主要参数}：标称阻值、精度（常见$\pm$5\%、$\pm$1\%）、功率等级、温度系数
\end{itemize}

电阻的功率计算公式：$P = I^2R = \frac{U^2}{R}$

上列各类电阻的特点可作如下补充说明。碳膜电阻成本低、噪声较大、温度系数
为负，多用于对精度要求不高的一般电路，精度等级常为 $\pm$5\%。金属膜电阻
噪声小、温度系数低、稳定性好，精度可达 $\pm$1\% 甚至更高，适用于测量
仪器与射频电路。线绕电阻功率大、精度高，适用于取样电阻与大功率负载。
热敏电阻的阻值随温度显著变化，分为负温度系数（NTC）与正温度系数（PTC）
两种，前者常用于测温与浪涌抑制，后者常用于过流保护与自恢复保险。压敏
电阻（MOV）在电压超过阈值时阻值急剧下降，用于吸收雷击与开关浪涌，是
通信设备电源入口防护的常用元件。

标称阻值并非任意取值，而是按 E 系列的几何级数分挡。精度 $\pm$5\% 的
电阻采用 E24 系列，每十倍频程分 24 挡；精度 $\pm$1\% 的电阻采用 E96
系列，每十倍频程分 96 挡。这样安排的目的，是使相邻挡位之间的间隔与
容差带宽相当，从而用最少的品种覆盖整个阻值范围。小型电阻的阻值多用
色环标注，四环电阻的前两环为有效数字、第三环为倍率、第四环为精度；
五环精密电阻的前三环为有效数字。片式电阻则直接用数字代码印字，如
“$103$” 表示 $10\times10^3\ \Omega = 10\ \mathrm{k}\Omega$。

多个电阻互连时的等效阻值遵循简单的规律：串联时 $R = R_1 + R_2 + \cdots$，
并联时 $\dfrac{1}{R} = \dfrac{1}{R_1} + \dfrac{1}{R_2} + \cdots$。
工程上常利用串并联组合得到非标称阻值，或用并联方式分摊功率。但需注意，
并联电阻的实际分流与各自阻值成反比，阻值偏差会导致功率分配不均，故用于
功率分摊时应选用同批次、同精度的电阻。

\subsection{电容器}

电容器用于存储电荷、滤波、耦合和调谐。
它由两块被介质隔开的导体极板构成，在直流稳态下相当于开路，因而具有
“隔直通交”的基本特性。正是这一特性，使电容器在电子设备中承担了四项
不可替代的职能：作为耦合元件在级间传送交流信号而隔断直流偏置；作为
旁路与滤波元件为交流分量提供低阻通路，平滑整流后的脉动直流；作为储能
元件在瞬态大电流时就近供电，抑制电源轨的跌落；作为谐振元件与电感构成
选频回路，决定振荡器与调谐放大器的工作频率。

平行板电容器的容量取决于介质性质与几何尺寸：
\[
C = \frac{\varepsilon_r \varepsilon_0 S}{d}
\]
其中 $\varepsilon_0$ 为真空介电常数，$\varepsilon_r$ 为介质的相对介电
常数，$S$ 为极板正对面积，$d$ 为极板间距。要在有限体积内获得大容量，
必须增大 $\varepsilon_r$ 与 $S$、减小 $d$。铝电解电容通过电化学方法在
铝箔表面生成极薄的氧化铝介质膜，使 $d$ 极小，因而单位体积容量很大；
陶瓷电容则利用高介电常数的陶瓷材料并采用多层叠合结构增大有效面积。

电容器所存储的电荷量与端电压成正比，即电容的充电公式为 $Q = CU$；
对时间求导即得其伏安关系 $i = C\,\mathrm{d}u/\mathrm{d}t$。这表明电容
两端电压不能突变，突变将导致无穷大电流。在电阻 $R$ 与电容 $C$ 串联并
接入恒压源 $U$ 的电路中，电容电压按指数规律上升：
\begin{equation}
u_C(t) = U\left(1 - e^{-t/RC}\right), \qquad i(t) = \frac{U}{R}e^{-t/RC}
\end{equation}
乘积 $\tau = RC$ 称为时间常数，经过 $\tau$ 时电容电压达到终值的约
$63.2\%$，经过 $5\tau$ 后已达终值的 $99\%$ 以上，工程上即视为充电完毕。
放电过程同理按 $u_C(t) = U_0 e^{-t/RC}$ 衰减。RC 时间常数是设计延时电路、
去耦网络与检波器负载的基本依据。

电容器储存的能量为 $W = \tfrac{1}{2}CU^2$。在正弦稳态下，其阻抗幅值为
$|Z_C| = 1/(\omega C) = 1/(2\pi f C)$，随频率升高而减小。但实际电容并
非理想元件：引线与极板存在等效串联电感 ESL，介质与电极存在等效串联电阻
ESR，三者串联构成一个串联谐振回路，其自谐振频率为
$f_0 = 1/(2\pi\sqrt{L_{\mathrm{ESL}}C})$。低于 $f_0$ 时呈电容性，高于
$f_0$ 时呈电感性，去耦作用随之失效。因此高频去耦必须选用引线短、ESL 小
的贴片电容，并常把大容量电解电容与小容量陶瓷电容并联使用，以在宽频带内
都获得低阻抗。介质损耗用损耗角正切 $\tan\delta = \omega C \cdot R_{\mathrm{ESR}}$
表示，损耗功率为 $P = I^2 R_{\mathrm{ESR}}$，这是开关电源滤波电容发热
乃至鼓包失效的主要原因。

\begin{itemize}
    \item \textbf{工作原理}：电容值由介电常数、极板面积和间距决定：$C = \frac{\varepsilon_r\varepsilon_0 S}{d}$
    \item \textbf{单位}：法拉（F），常用倍率：$\mu$F（微法）、nF（纳法）、pF（皮法）
    \item \textbf{类型}：电解电容、陶瓷电容、薄膜电容、云母电容、钽电容
    \item \textbf{主要参数}：标称容量、精度、额定电压、温度系数、等效串联电阻（ESR）
\end{itemize}

电容的充电公式：$Q = CU$

各类电容器的适用场合差别明显。铝电解电容容量大（可达数千 $\mu$F）、
价格低，但有极性、ESR 较大、寿命受温度制约，主要用于电源整流滤波与
低频耦合，反接会导致氧化膜击穿而爆裂。钽电解电容体积小、ESR 低、
频率特性优于铝电解，适用于小型化设备的电源去耦，但过压或反压时可能
短路起火，须严格降额。陶瓷电容分为温度稳定型（如 NP0/C0G）与高介电
常数型（如 X7R、Y5V）；前者容量小但温度系数极低、损耗极小，适合谐振
回路与定时电路，后者容量大但容量随温度和直流偏压明显下降，只宜用于
旁路与滤波。薄膜电容与云母电容的绝缘电阻高、损耗小、稳定性好，云母
电容尤其耐高压、耐大电流，传统上用于发射机的高频调谐与耦合回路。

\begin{table}[htbp]
\centering
\caption{常用电容器类型的特性对比}
\begin{tabular}{|p{2.0cm}|p{2.2cm}|p{2.4cm}|p{4.6cm}|}
\hline
\textbf{类型} & \textbf{容量范围} & \textbf{极性/损耗} & \textbf{典型用途与注意事项} \\
\hline
铝电解电容 & $1\ \mu$F$\sim$数千 $\mu$F & 有极性，ESR 大 & 电源整流滤波；严禁反接，高温下寿命显著缩短 \\
\hline
钽电容 & $0.1\sim470\ \mu$F & 有极性，ESR 小 & 小型设备电源去耦；须降额至额定电压的 $50\%$ 以下 \\
\hline
陶瓷电容 & 数 pF$\sim$数十 $\mu$F & 无极性，损耗小 & NP0/C0G 用于谐振定时，X7R 用于旁路去耦 \\
\hline
薄膜电容 & $1\ $nF$\sim$数十 $\mu$F & 无极性，损耗很小 & 音频耦合、交流滤波、吸收回路，稳定性好 \\
\hline
云母电容 & 数 pF$\sim$数 nF & 无极性，$Q$ 值高 & 高频高压调谐回路，精度与稳定性优良 \\
\hline
\end{tabular}
\end{table}

电容器的串并联规律与电阻相反：并联时容量相加 $C = C_1 + C_2 + \cdots$；
串联时 $\dfrac{1}{C} = \dfrac{1}{C_1} + \dfrac{1}{C_2} + \cdots$，
且各电容分压与容量成反比。串联可提高耐压，但因漏电流不均会造成分压
失衡，故实际电路中须在每只电容上并接均压电阻。可变电容与微调电容通过
改变极板正对面积或介质位置来调节容量，是传统调谐电路的核心部件。

\subsection{电感器}

电感器用于阻碍电流变化、滤波和储能。
它的本质是一段能够建立集中磁场的绕线：当流过的电流变化时，绕组产生的
自感电动势总是阻止该变化，故电感对交流呈现阻抗，对直流仅表现为绕线的
直流电阻。在通信设备中，电感与电容配合构成谐振回路以选择频率，构成
LC 滤波器以分离频段；作为扼流圈阻止射频电流窜入电源；在开关电源中
作为储能元件实现电压变换；绕在铁芯上互相耦合即成变压器，用于阻抗
变换与电气隔离。

对于长螺线管，其电感量为
\[
L = \frac{\mu N^2 S}{l}
\]
即与匝数 $N$ 的平方、磁芯磁导率 $\mu$、绕组截面积 $S$ 成正比，与磁路
长度 $l$ 成反比。可见增加匝数是提高电感量最有效的手段，而采用高磁导率
磁芯可在匝数不变的条件下大幅提高电感量。但磁芯存在饱和现象：当磁场强度
超过一定值后磁导率急剧下降，电感量随之减小，波形出现失真。因此电感的
额定电流不仅受绕线发热限制，还受磁芯饱和限制，两者取小者为准。

电感的伏安关系为 $u = L\,\mathrm{d}i/\mathrm{d}t$，说明流过电感的电流
不能突变。若强行切断电感电流，$\mathrm{d}i/\mathrm{d}t$ 极大将产生很高
的反向电动势，这就是继电器线圈、电机绕组断电时产生火花与浪涌的原因，
必须并接续流二极管或 RC 吸收网络加以抑制。电感的储能公式为
$W = \tfrac{1}{2}LI^2$，即能量与电流平方成正比而与电压无关，这与电容
恰成对偶。

在正弦稳态下电感阻抗为 $Z_L = \mathrm{j}\omega L$，其幅值 $|Z_L| = 2\pi f L$
随频率线性上升。由于绕线存在电阻 $R$，电感的品质因数定义为
\begin{equation}
Q = \frac{\omega L}{R} = \frac{2\pi f L}{R}
\end{equation}
$Q$ 值表征储能与损耗之比，直接决定谐振回路的选择性与插入损耗：
回路的通带带宽近似为 $B = f_0/Q$。$Q$ 值受多种损耗制约，包括绕线的
直流电阻、高频下的趋肤效应与邻近效应、磁芯的磁滞与涡流损耗、以及
屏蔽罩内的感生涡流。此外，绕组各匝之间存在分布电容 $C_0$，与电感
构成并联谐振，自谐振频率为
\[
f_{\mathrm{SRF}} = \frac{1}{2\pi\sqrt{L C_0}}
\]
超过该频率后元件呈电容性，不再具备电感作用。故高频电感应尽量减小
分布电容，常采用单层疏绕、分段绕制或蜂房式绕法。

\begin{itemize}
    \item \textbf{工作原理}：电感值与线圈匝数平方、磁芯磁导率和截面积成正比：$L = \frac{\mu N^2 S}{l}$
    \item \textbf{单位}：亨利（H），常用倍率：mH（毫亨）、$\mu$H（微亨）
    \item \textbf{类型}：空芯电感、铁氧体磁芯电感、铁粉芯电感、可变电感
    \item \textbf{主要参数}：标称电感量、精度、额定电流、品质因数（Q值）、分布电容
\end{itemize}

电感的储能公式：$W = \frac{1}{2}LI^2$

就上列类型而言，空芯电感不存在磁芯损耗与饱和问题，$Q$ 值高、线性好，
适用于甚高频与超高频回路，但电感量小、体积大且漏磁强。铁氧体磁芯电感
磁导率高、高频损耗小，是中短波至甚高频段扼流圈与滤波电感的主流选择；
按材料配方又分为镍锌系（适用于较高频率）与锰锌系（适用于较低频率、
磁导率更高）。铁粉芯（含铁硅铝等粉末芯）由绝缘包覆的磁性颗粒压制而成，
分布气隙使其抗饱和能力强、直流叠加特性好，广泛用于开关电源的储能电感
与差模滤波。可变电感通过旋动磁芯改变磁路耦合程度来微调电感量，用于
中频变压器与振荡回路的精确校准。

电感串联时（若无磁耦合）$L = L_1 + L_2 + \cdots$，并联时
$\dfrac{1}{L} = \dfrac{1}{L_1} + \dfrac{1}{L_2} + \cdots$；若两电感之间
存在磁耦合，则须计入互感 $M = k\sqrt{L_1L_2}$，其中 $k$ 为耦合系数。
同向串联时等效电感为 $L_1 + L_2 + 2M$，反向串联时为 $L_1 + L_2 - 2M$。
这一关系既是变压器与耦合谐振回路的设计基础，也提醒装配时相邻电感应
相互垂直放置或加以屏蔽，以避免意外耦合造成的自激振荡。

最后需要强调三种元件的协同关系。当电感与电容串联或并联时会出现谐振，
谐振频率为
\begin{equation}
f_0 = \frac{1}{2\pi\sqrt{LC}}
\end{equation}
串联谐振时回路阻抗最小、电流最大，可用于陷波与吸收；并联谐振时回路
阻抗最大，可用于选频与负载。回路中的电阻决定其 $Q$ 值与带宽：
$Q$ 高则选择性强但通带窄，$Q$ 低则通带宽而抑制差。整机调试中所谓
“调谐”，实质就是通过改变 $L$ 或 $C$ 使 $f_0$ 对准工作频率，并通过
适当的阻尼获得所需的带宽与增益平坦度。理解这三种无源元件的基本关系
与非理想特性，是分析一切通信电路的前提。
